\begin{figure}[t]
\centering
\includegraphics[clip, trim=0.75cm 0.3cm 0.3cm 0cm, width=\textwidth]{figures/architecture5.pdf} 
\caption{Architecture of the Step-Audio 2.}
\label{fig:architecture}
\end{figure}

\section{Methodology}
\subsection{Architecture}

Different from our previous Step-Audio~\cite{huang2025step},
Step-Audio 2 further integrates the generation of audio tokens into language modeling,
achieves end-to-end audio perception and generation.
As shown in Figure~\ref{fig:architecture}, Step-Audio 2 consists of
an audio encoder,
an audio adaptor,
an LLM decoder 
and an audio detokenizer.

The audio encoder is pretrained on various speech and audio understanding tasks including ASR, speaker age and gender prediction, audio event detection, etc.
The audio encoder has an output frame rate of 25 Hz and is frozen during the entire training process. 
An audio adaptor with a downsampling rate of 2 is employed to connect the audio encoder to LLM, thereby reducing the output frame rate of the audio encoder to 12.5 Hz. 

The LLM decoder directly takes the latent audio features from the audio adaptor as input,
and outputs an interleaved sequence of discrete text and audio tokens.
We employ the tokenizer from CosyVoice 2 \cite{du2024cosyvoice2} as the audio tokenizer.
And the text and audio tokens are interleaved \cite{huang2025stepaudioaqaa,huang2025step,zeng2024glm} at a fixed ratio and padded at the end to meet the ratio. %
The audio tokens are then extracted from the interleaved sequence and consumed by the audio detokenizer to generate the output waveform.
The input audio features and output interleaved sequences 
are then pre-filled as the history information for the next round of conversation.


To provide more accurate responses and expand interactive capabilities, %
we 
design tools to retrieve audio, current date and time, weather forecast and web content directly with explicit or implicit voice inputs.
Notably, 
we propose the audio search tool,
a novel tool with a voice library of hundreds of thousands of speeches with their corresponding transcriptions and descriptions.
With the retrieved speech from audio search, Step-Audio 2 is able to mimic the speaking style
or switching timbre according to the speech.
During inference,
the retrieved information is appended after the input audio features
before generating speech outputs.

Similar to Step-Audio~\cite{huang2025step} and Step-Audio-AQAA~\cite{huang2025stepaudioaqaa}
, Step-Audio 2’s audio detokenizer also consists of a Flow Matching module and a HiFi-GAN~\cite{kong2020hifigangenerativeadversarialnetworks} vocoder. The Flow-Matching module generates Mel spectrograms from the output audio tokens, while the vocoder further converts the Mel spectrograms into waveforms.
For Flow-Matching, we incorporate a CNN-based encoder layer after each self-attention module within the transformer block and train the model on 200,000 hours of high-quality speech. This enhancement significantly improves its Mel spectrogram reconstruction capability, leading to substantial gains in both pronunciation accuracy and timbre similarity.

Step-Audio 2 employs the same deployment infrastructure used in Step-Audio~\cite{huang2025step} and Step-Audio-AQAA~\cite{huang2025stepaudioaqaa}, which includes a voice activity detection (VAD) module to filter out input speeches and achieves real-time voice conversation. %







 







\subsection{Pre-training}

Step-Audio 2 model is initialized with a textual LLM and %
then continually pre-trained on 1.356T tokens of textual and audio data
over 21 days.

We first utilize
100B tokens
of ASR data to facilitate effective alignment between speech and text feature spaces within the adaptor.
During this phase, both the audio encoder and LLM are frozen, with only the adaptor being trained.
We conduct training for 12K steps at an 8,192 sequence length.
And the learning rate decays from $10^{-4}$ to $2\times10^{-5}$. 

We then extend the tokenizer of the textual LLM with 6.6K audio tokens.
To properly embed the new audio tokens and preserve the model's textual capabilities,
the model is then trained on 128B tokens of text data and 128B tokens of audio data.
Specifically, audio data includes 80B, 32B and 16B tokens of TTS, speech-to-speech conversation and utterance-level text-speech interleaved continuation data respectively.
The sequence length is increased to 16,384.
And the learning rates of the LLM,
adaptor, embedding layer
and output layer
are set to $2\times10^{-5}$, 
$5\times10^{-5}$, 
$5\times10^{-5}$, 
and $4\times10^{-5}$ respectively.

We then introduce our main pre-training process and further train the model on another 800B tokens of text and audio data.
We unify the learning rates to $2\times10^{-5}$ 
and
employ
400B tokens of textual data and 
42B, 120B, 8B, 30B, 5B, 45B and 150B tokens of ASR, TTS, speech-to-text translation, text-to-speech translation, speech-to-text continuation, utterance-level text-speech interleaved continuation and speech-to-speech conversation data respectively.

We finally employ 200B tokens of high-quality text and audio data to introduce a wider array of tasks and cooldown the model.
We employ
24.6B,
12.4B,
2.4B,
and 3.6B
tokens of audio data 
for 
multilingual and dialectal ASR, 
TTS, 
paralinguistic information understanding, 
speech-to-text translation respectively.
Besides, 
we develop a conversational speech synthesis pipeline to synthesize 6B, 15B and 36B tokens of audio data for speech-to-speech translation, utterance-level text-speech interleaved conversation and speech-to-speech conversation.
To ensure the vocal diversity in the synthesized speech, the system references a library of approximately 50k unique speakers.
We balance the audio data with 100B tokens of high-quality text data and
the learning rate decays from $2\times10^{-5}$ to $5\times10^{-6}$.

After this comprehensive pre-training procedure, the model 
has acquired strong audio understanding and generation capabilities
while maintaining its textual performance inherited from the initial textual LLM.































